\documentclass[11pt]{article}

\usepackage[T1]{fontenc}
\usepackage{amsmath,amssymb,amsthm,mathtools}
\usepackage[margin=1in]{geometry}
\usepackage[hidelinks]{hyperref}

\newcommand{\C}{\mathbb C}
\newcommand{\Z}{\mathbb Z}
\newcommand{\Sphere}{\widehat{\mathbb C}}
\newcommand{\Teich}{\operatorname{Teich}}
\newcommand{\PMod}{\operatorname{PMod}}
\newcommand{\Def}{\mathcal D}
\newcommand{\Par}{\mathcal P}
\newcommand{\id}{\operatorname{id}}
\newcommand{\im}{\operatorname{im}}
\newcommand{\wind}{\operatorname{wind}}

\newtheorem{theorem}{Theorem}[section]
\newtheorem{lemma}[theorem]{Lemma}
\newtheorem{proposition}[theorem]{Proposition}
\theoremstyle{definition}
\newtheorem{construction}[theorem]{Construction}
\theoremstyle{remark}
\newtheorem{remark}[theorem]{Remark}

\title{Problem 7 -- GPT-6 Astra}
\author{}
\date{}

\begin{document}
\maketitle

\begin{abstract}
We describe a simple closed curve for which the Dehn twist itself,
not merely a proper power, preserves the deformation space in the
question and has an infinite orbit on its connected components.
The curve is first specified for one explicit normalized rational map
and is then transported to an arbitrary map in the family.
The proof uses an explicit Hurwitz covering and a free-group quotient
of the stabilizer of a component.  It does not assume that the
components of the deformation space are simply connected.
\end{abstract}

\section{The explicit curve and the statement}

All marked points are labelled.  Thus mapping class groups below are
pure mapping class groups, and moduli spaces are moduli spaces of
ordered configurations.

\begin{lemma}[Normal form]\label{lem:normal}
After the unique M\"obius normalization
\[
c_1=0,\qquad f(c_1)=\infty,\qquad f^{\circ 2}(c_1)=1,
\]
the map in the question is
\begin{equation}\label{eq:normal}
F_r(z)=
1-\frac{r^2+r-1}{(r-1)z}
+\frac{r(2r-1)}{(r-1)z^2},
\end{equation}
with critical cycle
\[
0\longmapsto\infty\longmapsto1\longmapsto r\longmapsto0.
\]
Its other critical value is
\begin{equation}\label{eq:nu}
\nu(r)=
-\frac{(r^2-3r+1)^2}{4r(r-1)(2r-1)}.
\end{equation}
The allowed parameter space is
\begin{equation}\label{eq:parameter}
\Par=
\left\{
r\in\C:
r(r-1)(2r-1)
(r^2-3r+1)(r^2+r-1)(3r^2-3r+1)\ne0
\right\}.
\end{equation}
In particular, $\Par$ is path connected.
\end{lemma}

\begin{proof}
A quadratic rational map with $0$ critical, $F(0)=\infty$, and
$F(\infty)=1$ has the form
\[
F(z)=1+\frac b z+\frac c{z^2},\qquad c\ne0.
\]
The equations $F(1)=r$ and $F(r)=0$ give
\[
1+b+c=r,\qquad r^2+br+c=0,
\]
and solving them gives \eqref{eq:normal}.  The second critical value is
$1-b^2/(4c)$, which simplifies to \eqref{eq:nu}.
Writing $d(r)=4r(r-1)(2r-1)$, one also obtains
\[
\nu(r)-1=-\frac{(r^2+r-1)^2}{d(r)},\qquad
\nu(r)-r=-\frac{(3r^2-3r+1)^2}{d(r)}.
\]
These formulas give precisely \eqref{eq:parameter}.  Its complement
in $\C$ is finite, so $\Par$ is path connected.
\end{proof}

Set
\[
R=100,\qquad
v_*=\nu(R)=-\frac{9701^2}{4\cdot100\cdot99\cdot199}.
\]
Notice that
\begin{equation}\label{eq:vbound}
-12<v_*<-11.
\end{equation}
Write
\[
A_*=\{0,\infty,1,R\},\qquad B_*=A_*\cup\{v_*\}.
\]
For a simple closed curve $a$, let $T_a$ denote the positive Dehn twist
about $a$.

\begin{construction}[The model curve]\label{con:model}
On $(\Sphere,B_*)$, take the circle
\[
C=\{z:|z-50|=55\}.
\]
Let $\alpha$ be the embedded polygonal arc
\[
[v_*,v_*-200i]\ \cup\
[v_*-200i,R-200i]\ \cup\
[R-200i,R].
\]
Let $\delta$ be the boundary of a sufficiently thin regular
neighborhood of $\alpha$, for example its $1/10$-neighborhood with
corners rounded.  Define
\begin{equation}\label{eq:modelcurve}
\boxed{\ \gamma_*=T_\delta(C).\ }
\end{equation}
Here any representative of $T_\delta$ supported in a thin collar of
$\delta$ may be used; the resulting isotopy class of $\gamma_*$
relative to $B_*$ is independent of that choice.
\end{construction}

Thus $\delta$ surrounds precisely $v_*$ and $R$, whereas $C$, and hence
$\gamma_*$, separates $\{v_*,\infty\}$ from $\{0,1,R\}$.
The twist in \eqref{eq:modelcurve} is important: the partition of the
marked points alone does not specify the required curve.

Here is the corresponding explicit prescription for the map in the
question.

\begin{construction}[Transport to an arbitrary map]\label{con:general}
Let $m$ be the normalization in Lemma~\ref{lem:normal}, so that
$mfm^{-1}=F_{r_0}$ for some $r_0\in\Par$.

Choose a piecewise smooth path $r(t)$ in $\Par$ from $r_0$ to $R$.
This choice can be made entirely explicit: take the polygonal path
\[
r_0\longrightarrow N+iN^2\longrightarrow R,
\]
where $N$ is the least positive integer for which both segments avoid
$\C\setminus\Par$.  Such an $N$ exists, because the forbidden
intermediate points lie on finitely many real lines, each meeting
the parabola $\{x+ix^2:x\in\mathbb R\}$ finitely often.

Choose an ambient isotopy $h_t$ of $\Sphere$, starting at the identity,
such that
\[
h_t(0)=0,\quad h_t(1)=1,\quad h_t(\infty)=\infty,\quad
h_t(r_0)=r(t),\quad h_t(\nu(r_0))=\nu(r(t)).
\]
Define
\begin{equation}\label{eq:generalcurve}
\boxed{\ \gamma_f=m^{-1}\bigl(h_1^{-1}(\gamma_*)\bigr).\ }
\end{equation}
Existence of $h_t$, and independence of its choice up to the relevant
isotopy, are justified in Section~\ref{sec:transport}.
\end{construction}

\begin{theorem}\label{thm:main}
The curve $\gamma_f$ of Construction~\ref{con:general} satisfies
\[
T_{\gamma_f}(\Def_f)=\Def_f.
\]
If $\mathcal C_f$ is the connected component containing the standard
complex structure of the given rational map, then
\[
T_{\gamma_f}^{\,n}(\mathcal C_f),\qquad n\in\Z,
\]
are pairwise distinct connected components of $\Def_f$.
Consequently, the power requested in the question can be taken to be
$1$.
\end{theorem}

The Hurwitz-space coordinates used below are those of
Hironaka--Koch \cite[Proposition~3.1]{HK}; we verify the required
properties explicitly.  Related pairs of equalizing curves occur in
Firsova--Kahn--Selinger \cite[Section~9, case~(12b), arXiv version]{FKS}.
The argument below identifies a particular single twist and proves
its infinite component orbit.
% Sources for the cited coordinate model and related curve types:
% ([math.fsu.edu](https://www.math.fsu.edu/~hironaka/papers/disconnected.pdf))

\section{A covering-space model for the deformation space}

Put
\[
M=\C\setminus\{0,1\},\qquad
Y=\{(y,z)\in M^2:y\ne z\}.
\]
These are the moduli spaces of the configurations
$(0,\infty,1,x)$ and $(0,\infty,1,y,z)$, respectively.

We use the standard facts that the Teichm\"uller spaces of these
labelled punctured spheres are their universal covers, with deck
groups their pure mapping class groups.  Contractibility and the
properly discontinuous action are standard Teichm\"uller theory;
the labelled genus-zero action is free because a M\"obius map fixing
three labelled points is the identity.  See
\cite[Chapters~10 and~12]{FM}, or the formulation for marked spheres in
\cite[Section~2.2]{HK}.
% Source for these standard Teichmuller-covering facts:
% ([pagine.dm.unipi.it](https://pagine.dm.unipi.it/~a019210/Farb%20Magalit_Primer%20on%20Teichmuller%20theory.pdf))

We let a mapping class $h$ act on markings by
\[
h\cdot[\varphi]=[\varphi\circ h^{-1}].
\]
Accordingly, a configuration loop extended to an ambient isotopy
$H_t$ determines the deck element $[H_1^{-1}]$.
Peripheral loops will always be oriented to represent the indicated
Dehn twist.

\subsection{The explicit Hurwitz covering}

For $(x,y)\in\C^2$, define
\[
\begin{split}
D(x,y)&=x(x-1)(x+y-1),\\
Q_0(x,y)&=(x-1)^2-y,\\
Q_1(x,y)&=x^2+y-1,\\
Q_y(x,y)&=(x-1)^2+(2x-1)y,
\end{split}
\]
and
\begin{equation}\label{eq:Z}
Z(x,y)=-\frac{Q_0(x,y)^2}{4D(x,y)}.
\end{equation}
Let
\begin{equation}\label{eq:W}
W=
\left\{(x,y):
y(y-1)D(x,y)Q_0(x,y)Q_1(x,y)Q_y(x,y)\ne0
\right\}.
\end{equation}

\begin{lemma}\label{lem:cover}
The map
\[
q:W\longrightarrow Y,\qquad q(x,y)=(y,Z(x,y)),
\]
is a connected, unbranched covering of degree four.
The two maps from $W$ to $M$ recording the source configuration and
the target configuration with the extra critical value forgotten are
\[
p_x(x,y)=x,\qquad p_y(x,y)=y.
\]
\end{lemma}

\begin{proof}
A normalized quadratic map satisfying
\[
0\longmapsto\infty,\qquad
\infty\longmapsto1,\qquad
1\longmapsto y,\qquad x\longmapsto0
\]
must be
\begin{equation}\label{eq:Fxy}
G_{x,y}(t)=\frac{(t-x)(t-s)}{t^2},
\qquad s=\frac{x+y-1}{x-1}.
\end{equation}
Indeed, the first two conditions make the numerator monic, and
$G_{x,y}(1)=y$ determines $s$.
Its second critical value is
\[
1-\frac{(x+s)^2}{4xs}=Z(x,y).
\]
Direct expansion gives
\begin{equation}\label{eq:Zidentities}
Z-1=-\frac{Q_1^2}{4D},\qquad
Z-y=-\frac{Q_y^2}{4D},\qquad
\frac{\partial Z}{\partial x}
=-\frac{Q_0Q_1Q_y}{4D^2}.
\end{equation}
Thus \eqref{eq:W} says exactly that the source and target markings
are distinct and that the second critical value avoids the other
target markings.

For fixed $(y,z)\in Y$, the equation $Z(x,y)=z$ is
\[
Q_0(x,y)^2+4zD(x,y)=0,
\]
a monic polynomial equation of degree four in $x$.
None of its roots is $0$, $1$, or $1-y$:
at these points $Q_0$ is nonzero because $y\ne0,1$.
The identities \eqref{eq:Zidentities} show that all four roots are
simple when $z\notin\{0,1,y\}$.  The implicit function theorem gives
four local inverse branches of $q$.  Hence $q$ is an unbranched
four-sheeted covering.

Finally, $W$ is path connected.  Given two of its points, restrict
the nonvanishing polynomial defining $W$ to the complex affine line
through them.  It is not identically zero, and therefore excludes only
finitely many points of that line.  The two points can be joined in
the complement of those finitely many points.

The descriptions of $p_x$ and $p_y$ follow from
\eqref{eq:Fxy}.
\end{proof}

\begin{lemma}\label{lem:teich}
For each $r\in\Par$, there is a universal covering
\[
U_r:\Teich(B_r)\longrightarrow W,\qquad
B_r=\{0,\infty,1,r,\nu(r)\},
\]
taking the standard marked structure $\tau_r$ to $(r,r)$, such that
\begin{equation}\label{eq:coverdiagram}
qU_r=\pi_{B_r},\qquad
\pi_{A_r}F_r^*=p_xU_r,\qquad
\pi_{A_r}i_*=p_yU_r,
\end{equation}
where $A_r=\{0,\infty,1,r\}$.
\end{lemma}

\begin{proof}
Given a marked complex structure $[\varphi]\in\Teich(B_r)$, normalize
$\varphi$ to fix $0,\infty,1$.  Uniformize the pullback complex
structure, obtaining a normalized marking $\psi$ for which
\[
G=\varphi\circ F_r\circ\psi^{-1}
\]
is rational.  This is the usual uniformization description of the
pullback map; see \cite[Sections~2.1--2.2, Proposition~2.1]{HK}.
Set
\[
x=\psi(r),\qquad y=\varphi(r),\qquad z=\varphi(\nu(r)).
\]
The map $G$ has precisely the normalization in \eqref{eq:Fxy}, so
$z=Z(x,y)$.  This gives the continuous lift $U_r$ and all identities
in \eqref{eq:coverdiagram}.

Since $\Teich(B_r)$ is the universal cover of $Y$ and $q$ is a
connected covering, its lift $U_r$ is the universal cover of $W$.
This is the standard composition and lifting property of coverings;
see \cite[Section~1.3]{H}.
\end{proof}

\subsection{The stabilizer of one component}

For the rest of the model argument, use $r=R$, write $U=U_R$, and
take the basepoints
\[
\tau_*=\tau_R,\qquad w_*=(R,R).
\]
The injection $q_*$ identifies
\[
L:=\pi_1(W,w_*)
\]
with a subgroup of the deck group of $\Teich(B_*)\to Y$.

Define
\[
d:\Par\longrightarrow W,\qquad d(r)=(r,r),
\qquad
K:=\im\bigl(d_*:\pi_1(\Par,R)\to L\bigr).
\]
The formulas defining $W$ show that its diagonal is exactly
$d(\Par)$.

\begin{lemma}\label{lem:stabilizer}
Let $\mathcal C_*$ be the component of $\Def_{F_R}$ containing
$\tau_*$.  Then:
\begin{enumerate}
\item $\Def_{F_R}$ is a union of connected components of
$U^{-1}(d(\Par))$, and each such component is path connected;
\item for $g\in L$,
\[
g(\mathcal C_*)=\mathcal C_*
\quad\Longleftrightarrow\quad g\in K;
\]
\item if
\[
(p_x)_*(g)=(p_y)_*(g)
\quad\text{in }\pi_1(M,R),
\]
then $g(\Def_{F_R})=\Def_{F_R}$.
\end{enumerate}
\end{lemma}

\begin{proof}
On $U^{-1}(d(\Par))$, the maps $F_R^*$ and $i_*$ are lifts to
$\Teich(A_*)$ of the same map to $M$, by
\eqref{eq:coverdiagram}.  On a connected component, two lifts either
agree everywhere or nowhere: if they agree at one point, uniqueness
of lifting makes them agree throughout that component.
Since $U^{-1}(d(\Par))\to d(\Par)$ is a covering of a Riemann surface,
its components are path connected.  This proves the first assertion.

A deck element $g\in L$ carries $\tau_*$ into $\mathcal C_*$ exactly
when there is a path in $U^{-1}(d(\Par))$ from $\tau_*$ to
$g(\tau_*)$.  Projecting this path gives a loop in $d(\Par)$
representing $g$; conversely, such a loop lifts to such a path.
This proves the second assertion.

For the last assertion, lifting uniqueness gives the equivariance
relations
\[
F_R^*\circ g=(p_x)_*(g)\circ F_R^*,
\qquad
i_*\circ g=(p_y)_*(g)\circ i_*,
\]
where the elements on the right act as deck transformations of
$\Teich(A_*)$.  Equality of the two deck elements implies
$g(\Def_{F_R})\subseteq\Def_{F_R}$; applying the same argument to
$g^{-1}$ gives equality.
\end{proof}

This is the covering-space form of the special-liftable and
component-stabilizer descriptions in
\cite[Sections~2.8--2.9]{HK}.  The covering facts used in its proof are
the lifting and deck-transformation results in
\cite[Section~1.3, especially Proposition~1.39]{H}.
% Sources for the cited stabilizer and covering-space descriptions:
% ([math.fsu.edu](https://www.math.fsu.edu/~hironaka/papers/disconnected.pdf))

\section{A free-group quotient detecting movement of components}

Enlarge $W$ to the five-line complement
\[
X=\{(x,y)\in M^2:x+y\ne1\},
\]
and let $j:W\hookrightarrow X$ be inclusion.  Put
\[
\Par_0=\C\setminus\{0,1,\tfrac12\},\qquad
e:\Par_0\longrightarrow X,\quad e(r)=(r,r).
\]
This auxiliary complement also appears in
\cite[Section~2.12]{HK}.  The properties needed here have short direct
proofs.
% Source for the auxiliary five-line complement:
% ([math.fsu.edu](https://www.math.fsu.edu/~hironaka/papers/disconnected.pdf))

\begin{lemma}\label{lem:freeimage}
The map $e_*$ is injective.  Consequently
\[
H:=e_*\pi_1(\Par_0,R)
\]
is a free group of rank three.  Moreover,
\begin{equation}\label{eq:imageK}
j_*(K)=H.
\end{equation}
\end{lemma}

\begin{proof}
Consider
\[
\Phi:X\longrightarrow M,\qquad
\Phi(x,y)=\frac{xy}{x+y-1}.
\]
It avoids $0$ because $xy\ne0$, and it avoids $1$ because
\[
xy-(x+y-1)=(x-1)(y-1)\ne0.
\]
On the diagonal,
\[
\Phi e(r)=\psi(r):=\frac{r^2}{2r-1}.
\]
The map $\psi:\Par_0\to M$ is an unbranched double covering.
Indeed, for $w\notin\{0,1\}$ its inverse points are the two distinct
roots of
\[
r^2-2wr+w=0,
\]
and neither root is $0$, $1$, or $1/2$.
Thus $\psi_*$ is injective, as is the fundamental-group homomorphism
of any covering.  Since $\Phi_*e_*=\psi_*$, $e_*$ is injective.

The plane with three points removed has fundamental group free of
rank three: disjoint small circles about the punctures, joined to a
basepoint by a tree of arcs, give the usual graph model.
Thus $H\cong F_3$.

Finally, $\Par$ is obtained from $\Par_0$ by deleting finitely many
additional points.  The inclusion $\Par\hookrightarrow\Par_0$
induces a surjection on fundamental groups: a loop can be perturbed
in small disks to avoid those additional points.
Using $jd=e|_{\Par}$ proves \eqref{eq:imageK}.
\end{proof}

We will also use the winding-number homomorphism
\begin{equation}\label{eq:lambda}
\lambda:\pi_1(X,w_*)\longrightarrow\Z,\qquad
\lambda([\beta])=\wind_0\bigl((x+y-1)\circ\beta\bigr).
\end{equation}

\begin{lemma}[A free-group centralizer criterion]\label{lem:centralizer}
Let $F$ be a free group and $\theta:F\to G$ a homomorphism.
If $\theta(t)$ has infinite order, then
\[
\{a\in F:at=ta,\ \theta(a)=1\}=\{1\}.
\]
\end{lemma}

\begin{proof}
The centralizer of a nonidentity element of a free group is infinite
cyclic.  To recall the argument, let the free group act on its Cayley
tree.  Writing an element as a conjugate of a cyclically reduced word
exhibits its translation axis, which is unique.  A commuting element
preserves this axis and acts on it by a translation, not a reflection.
Signed translation distance embeds the centralizer into $\Z$:
an element with zero translation fixes a vertex, and the left action
of a free group on its Cayley tree is free on vertices.
The centralizer is therefore cyclic.

Write this centralizer as $\langle b\rangle$, with $t=b^k$, $k\ne0$.
Since $\theta(t)$ has infinite order, so does $\theta(b)$.
Hence $\theta$ is injective on $\langle b\rangle$, proving the claim.
\end{proof}

\section{Verification for the explicit model curve}

\subsection{Elementary twist conventions}

We will use three elementary facts:
\[
T_{h(a)}=hT_ah^{-1},
\]
twists about disjoint curves commute, and a twist about the boundary
of a disk containing exactly two marked points is represented by
moving one of those points once around the other.

For completeness, the first two facts follow by transporting, or
choosing disjoint, annular supports.  For the third, identify the disk
with a round disk, place the fixed puncture at its center, and use
\[
z\longmapsto e^{2\pi it\chi(|z|)}z,
\]
where $\chi$ is one on a disk containing the moving puncture and zero
near the boundary.  The puncture travels once around the center, and
the endpoint is one boundary Dehn twist, with sign determined by the
chosen conventions.  Reversing the loop gives the other sign.
These are also the standard twist and point-pushing facts in
\cite[Sections~3.3 and~4.2]{FM}.
% Source for the cited twist and point-pushing facts:
% ([pagine.dm.unipi.it](https://pagine.dm.unipi.it/~a019210/Farb%20Magalit_Primer%20on%20Teichmuller%20theory.pdf))

Let
\[
\eta=\{z:|z|=2\},\qquad
u=T_\delta,\qquad
g=T_{\gamma_*}=uT_Cu^{-1}.
\]
The disk bounded by $\eta$ lies strictly inside the disk bounded by
$C$, and the neighborhood defining $\delta$ is disjoint from that
disk.  Therefore
\begin{equation}\label{eq:disjoint}
\gamma_*\cap\eta=\varnothing
\end{equation}
for suitable representatives.

\subsection{The monodromy of the required twist}

\begin{lemma}\label{lem:meridian}
The element $g=T_{\gamma_*}$ belongs to $L$ and satisfies
\begin{equation}\label{eq:gproperties}
(p_x)_*(g)=1,\qquad
(p_y)_*(g)=1,\qquad
|\lambda(j_*g)|=1.
\end{equation}
In particular,
\[
g(\Def_{F_R})=\Def_{F_R}.
\]
\end{lemma}

\begin{proof}
Keep $y=R$ fixed and consider the rational map
\[
x\longmapsto Z(x,R).
\]
By \eqref{eq:Zidentities}, its branch values are $0,1,R$.
Each has two simple ramification points.  Its four poles
\[
0,\quad1,\quad1-R,\quad\infty
\]
are simple.

The M\"obius involution
\begin{equation}\label{eq:involution}
J(x)=\frac{(1-R)(x-1)}{x+R-1}
\end{equation}
satisfies
\begin{equation}\label{eq:Jidentities}
J^2=\id,\qquad Z(J(x),R)=Z(x,R).
\end{equation}
Both identities follow by substitution.  Its fixed-point equation is
\[
x^2+2(R-1)x+1-R=0,
\]
which is exactly $Q_y(x,R)=0$.
Thus its fixed points are the two ramification points over $R$.

It follows that the sheet monodromy of a based meridian about the
critical value $R$ is $J$.  Indeed, close to either of these
ramification points a meridian interchanges the two local inverse
branches, as does $J$.  Transport to the basepoint commutes with
$J$ by \eqref{eq:Jidentities}.  The assertion holds for either
orientation of the meridian.

By the point-pushing description, $u=T_\delta$ is represented by such
a meridian about $z=R$, with $y=R$ fixed.  Consequently its lift to
the four-sheeted covering starts at $x=R$ and ends at $x=J(R)$.

The twist $T_C$ is represented by a meridian about $z=\infty$.
Its access path from $v_*$ to infinity may be chosen along the negative
real axis, since that ray lies in the complementary disk of $C$
containing $v_*$ and $\infty$.

Starting on the sheet $x=R$, the lift of this access path tends to
$x=\infty$.  To see this directly, for real $x\ge R$ all three
quantities $Q_0(x,R),Q_1(x,R),Q_y(x,R)$ are positive.
Equation \eqref{eq:Zidentities} therefore gives
\[
\frac{d}{dx}Z(x,R)<0.
\]
Moreover,
\[
Z(R,R)=v_*,\qquad Z(x,R)\longrightarrow-\infty
\quad(x\longrightarrow+\infty).
\]
Thus this real inverse branch parametrizes precisely the required
lift.

Starting instead at $J(R)$, the lifted access path is its image under
$J$, and therefore tends to
\[
J(\infty)=1-R.
\]
Because this is a simple pole, the infinity meridian lifts to a
single meridian about $x=1-R$, and returns to the same sheet.

Now represent
\[
g=uT_Cu^{-1}
\]
by the corresponding conjugated configuration loop.  Its lift first
changes the sheet from $R$ to $J(R)$, then makes the just-described
meridian about $x=1-R$, and finally returns along the inverse
conjugating path.  Hence the lift is closed, proving $g\in L$.
In the fiber $y=R$, its homotopy class is a based conjugate of a
single meridian about $x=1-R$.

Since $1-R\notin\{0,1\}$, that meridian is null-homotopic after
projection to the $x$-coordinate space $M$.  Its $y$-coordinate is
constant.  On the other hand, $x+y-1=x+R-1$ winds once, with one of
the two signs, about zero.  This proves \eqref{eq:gproperties}.

The final assertion follows from Lemma~\ref{lem:stabilizer}(3).
\end{proof}

\subsection{A commuting twist that does preserve the base component}

\begin{lemma}\label{lem:cusp}
The diagonal loop
\begin{equation}\label{eq:tloop}
t(s)=\bigl(Re^{2\pi is},Re^{2\pi is}\bigr),
\qquad 0\le s\le1,
\end{equation}
defines an element $t\in K$ whose mapping class is
$T_\eta$ or $T_\eta^{-1}$.
Moreover, $(p_y)_*(j_*t)$ has infinite order in $\pi_1(M,R)$.
\end{lemma}

\begin{proof}
The defining factors of \eqref{eq:parameter} do not vanish on
$|r|=R=100$, by the triangle inequality.  Thus \eqref{eq:tloop}
is a loop in $d(\Par)$, so $t\in K$.

We identify its mapping class.  A calculation from \eqref{eq:nu}
gives
\[
\frac{\nu(r)}r+\frac18
=
\frac{9r^3-21r^2+12r-2}
     {8r^2(r-1)(2r-1)}.
\]
Consequently, when $|r|=R=100$,
\begin{equation}\label{eq:estimate}
\left|\frac{\nu(r)}r+\frac18\right|
\le
\frac{9R^3+21R^2+12R+2}
     {8R^2(R-1)(2R-1)}
<\frac1{100}.
\end{equation}
Set $a_0=\nu(R)/R$.
The projected configuration loop
\[
\bigl(r(s),\nu(r(s))\bigr),\qquad r(s)=Re^{2\pi is},
\]
has the based homotopy
\[
\left(
r(s),\
r(s)\left((1-a)\frac{\nu(r(s))}{r(s)}+aa_0\right)
\right),
\qquad 0\le a\le1.
\]
Indeed, the expression in parentheses stays in the convex disk
of radius $1/100$ about $-1/8$.
The moving critical value therefore has modulus greater than
$11.5$, and is never $0$, $1$, or $r(s)$.
At the end of the homotopy the loop is
\[
\bigl(Re^{2\pi is},v_*e^{2\pi is}\bigr).
\]

Choose a smooth radial function $\chi$ equal to zero on $[0,2]$
and equal to one on $[10,\infty)$.
The ambient isotopy
\[
H_s(z)=e^{2\pi is\chi(|z|)}z,\qquad H_s(\infty)=\infty,
\]
realizes this last configuration loop.  Its endpoint is one Dehn
twist in the unmarked annulus $2<|z|<10$, and hence is a twist
about $\eta$.  Passing to the deck-action convention may invert
it, which proves the asserted description of $t$.

Finally, the $y$-coordinate of \eqref{eq:tloop} winds once about zero.
Its image in $\pi_1(M,R)$ therefore has infinite order.
\end{proof}

\subsection{Infinitely many different components}

\begin{proposition}\label{prop:model}
The components
\[
g^n(\mathcal C_*),\qquad n\in\Z,
\]
are pairwise distinct.
\end{proposition}

\begin{proof}
By Lemma~\ref{lem:meridian}, $g$ preserves the deformation space.
Suppose that
\[
g^n(\mathcal C_*)=\mathcal C_*
\]
for some nonzero integer $n$.  Lemma~\ref{lem:stabilizer}(2) gives
$g^n\in K$, and hence, by Lemma~\ref{lem:freeimage},
\[
j_*(g^n)\in H.
\]
Also $j_*t\in H$.

The curves $\gamma_*$ and $\eta$ are disjoint by
\eqref{eq:disjoint}.  Thus $g$ commutes with the mapping class of
$t$.  Since $q_*$ is injective, they commute in $L$, and their images
commute in $\pi_1(X,w_*)$.

Apply Lemma~\ref{lem:centralizer} to the free group $H$ and the
homomorphism
\[
(p_y)_*|_H:H\longrightarrow\pi_1(M,R).
\]
By Lemma~\ref{lem:cusp}, the image of $j_*t$ has infinite order.
By Lemma~\ref{lem:meridian}, the image of $j_*(g^n)$ is the identity.
The centralizer criterion therefore gives
\[
j_*(g^n)=1.
\]
But \eqref{eq:gproperties} and \eqref{eq:lambda} imply
\[
\left|\lambda\bigl(j_*(g^n)\bigr)\right|=|n|>0,
\]
a contradiction.

Thus no nonzero power of $g$ stabilizes $\mathcal C_*$.
If two of the displayed components were equal, their index
difference would give such a power.  They are therefore pairwise
distinct.
\end{proof}

\section{Transport and completion of the proof}\label{sec:transport}

We first justify the isotopy required in
Construction~\ref{con:general}.
Along the compact parameter path, the five labelled points
\[
0,\quad1,\quad\infty,\quad r(t),\quad\nu(r(t))
\]
remain distinct.  On each smooth portion of the path, choose disjoint
small disks around the two moving finite points, avoiding $0$ and $1$.
A compactly supported time-dependent vector field which equals the
velocity of each moving point near that point has a flow tracking
both points and fixing $0,1,\infty$.  Concatenating these flows
constructs $h_t$.

If $h_t$ and $\widetilde h_t$ are two such extensions of the same
configuration path, then
\[
\widetilde h_t^{-1}h_t
\]
fixes the initial labelled configuration pointwise for every $t$.
Thus the two endpoint pullbacks of $\gamma_*$ are isotopic relative
to that configuration.  This proves that the construction specifies
an isotopy class without any further choice of ambient isotopy.

\begin{lemma}\label{lem:transport}
For the isotopy $h_t$ of Construction~\ref{con:general}, put $h=h_1$.
There is a homeomorphism
\[
\Theta:\Teich(B_*)\longrightarrow\Teich(B_{r_0}),
\qquad
\Theta([\chi])=[\chi\circ h],
\]
such that
\[
\Theta(\Def_{F_R})=\Def_{F_{r_0}}.
\]
It takes $\mathcal C_*$ onto the component containing the standard
complex structure for $F_{r_0}$, and intertwines $T_{\gamma_*}$ with
$T_{h^{-1}(\gamma_*)}$.
\end{lemma}

\begin{proof}
Consider the path
\[
\tau(t)=[h_t]\in\Teich(B_{r_0}).
\]
Its projection to $Y$ is
\[
\bigl(r(t),\nu(r(t))\bigr).
\]
The diagonal path $(r(t),r(t))$ is a lift of this path to $W$.
By Lemma~\ref{lem:teich} and uniqueness of lifting,
\[
U_{r_0}(\tau(t))=(r(t),r(t)).
\]
The two paths $F_{r_0}^*(\tau(t))$ and $i_*(\tau(t))$ are therefore
lifts to $\Teich(A_{r_0})$ of the same path $r(t)$ in $M$.
They start at the same point, so they coincide for every $t$.
Consequently
\[
\tau(t)\in\Def_{F_{r_0}}.
\]

At $t=1$, normalized uniformization of the pullback gives a marking
$h'$ such that
\begin{equation}\label{eq:equivalence}
F_R\circ h'=h\circ F_{r_0}.
\end{equation}
The rational map on the left is $F_R$ because its Hurwitz coordinates
are $(R,R)$, and \eqref{eq:Fxy} uniquely determines that normalized
map.  Since $\tau(1)$ lies in the deformation space,
\[
h'\simeq h\quad\text{relative to }A_{r_0}.
\]

Let
\[
\Theta_A:\Teich(A_*)\longrightarrow\Teich(A_{r_0})
\]
be change of marking by $h|_{A_{r_0}}$.
Naturality of pullback applied to \eqref{eq:equivalence}, and the
relative isotopy of $h'$ and $h$, give
\[
F_{r_0}^*\Theta=\Theta_A F_R^*,
\qquad
i_*\Theta=\Theta_A i_*.
\]
Since $\Theta_A$ is a homeomorphism, these identities imply
\[
\Theta(\Def_{F_R})=\Def_{F_{r_0}}.
\]

Furthermore, $\Theta(\tau_*)=[h]=\tau(1)$ lies in the same component
as the standard structure $\tau(0)$.  Hence $\Theta$ carries
$\mathcal C_*$ onto that component.

Finally, change of marking conjugates a mapping class $a$ to
$h^{-1}ah$.  Since conjugating a Dehn twist gives the Dehn twist
about the inverse-image curve, $\Theta$ intertwines the stated
twists.
\end{proof}

\begin{proof}[Proof of Theorem~\ref{thm:main}]
Lemma~\ref{lem:meridian} and Proposition~\ref{prop:model} prove the
theorem for $F_R$ and $\gamma_*$.
Lemma~\ref{lem:transport} transfers both preservation of the
deformation space and pairwise distinctness of the component orbit
to $F_{r_0}$ and $h^{-1}(\gamma_*)$.
Undoing the M\"obius normalization transfers these conclusions to
$f$ and the curve \eqref{eq:generalcurve}.
\end{proof}

\begin{remark}
The proof distinguishes two assertions which should not be conflated.
The vanishing of both coordinate monodromies proves that the twist
preserves the \emph{whole} deformation space.
Its failure to preserve any nonzero-power iterate of the base
component is a separate fact, detected by the winding homomorphism
\eqref{eq:lambda} together with a cyclic-centralizer argument in the
free group $H$.
\end{remark}

\begin{thebibliography}{9}

\bibitem{FM}
B.~Farb and D.~Margalit,
\emph{A Primer on Mapping Class Groups},
Princeton Mathematical Series, vol.~49,
Princeton University Press, Princeton, NJ, 2012.

\bibitem{FKS}
T.~Firsova, J.~Kahn, and N.~Selinger,
On deformation spaces of quadratic rational functions,
\emph{International Mathematics Research Notices}
\textbf{2023} (2023), no.~8, 6703--6738.
doi:10.1093/imrn/rnab320.
The arXiv version referred to above is
arXiv:1606.00723v4.
% Bibliographic verification:
% ([academic.oup.com](https://academic.oup.com/imrn/article-abstract/2023/8/6703/6545985))

\bibitem{H}
A.~Hatcher,
\emph{Algebraic Topology},
Cambridge University Press, Cambridge, 2002.

\bibitem{HK}
E.~Hironaka and S.~Koch,
A disconnected deformation space of rational maps,
\emph{Journal of Modern Dynamics}
\textbf{11} (2017), 409--423.
doi:10.3934/jmd.2017016.
% Bibliographic verification:
% ([aimsciences.org](https://www.aimsciences.org/article/doi/10.3934/jmd.2017016))

\end{thebibliography}

\end{document}